\documentclass[11pt]{article}
% Shared preamble for the round 1 response letters.
% RED-D-26-00078, submitted as "Housing First or Treatment First?",
% resubmitted as "Housing First or Treatment First? Ask the Kids, Not the Parents"
%
% House style, decided once here:
%   * responses are in RED, referee text is in the black italic block above them
%   * NO HYPHENS anywhere in the prose
%   * every response ends with a pointer to where in the new paper it lives

\usepackage[margin=1.05in]{geometry}
\usepackage{amsmath,amssymb}
\usepackage{xcolor}
\usepackage[official]{eurosym}
\usepackage{booktabs}
\usepackage{longtable}
\usepackage{array}
\usepackage[hidelinks]{hyperref}

% enumitem is not in this TeX installation, so first level list spacing is set
% directly instead
\makeatletter
\renewcommand{\@listI}{%
  \leftmargin=1.7em \labelwidth=1.2em
  \topsep=0.4em \parsep=0pt \itemsep=0.3em \partopsep=0pt}
\makeatother

\definecolor{respred}{RGB}{190,25,25}

\setlength{\parskip}{0.35em}
\setlength{\parindent}{0pt}

\newcolumntype{L}[1]{>{\raggedright\arraybackslash}p{#1}}
\newcommand{\pname}[1]{\texttt{#1}}

% ---- the two building blocks -------------------------------------------------
% \cmt{label}{short title} opens a comment, then the verbatim or paraphrased
% referee text goes in the referee environment, then the response.

\newcommand{\cmt}[2]{%
  \par\medskip
  \noindent\textbf{#1\quad #2}\par}

\newenvironment{referee}
  {\begin{list}{}{\leftmargin=1.6em \rightmargin=1.2em \topsep=0.4em}\item[]\small\itshape}
  {\end{list}}

\newenvironment{response}
  {\par\begingroup\color{respred}\medskip}
  {\par\endgroup\medskip}

% a flag for statements in the reports that are factually wrong, so they are
% easy to find again and easy to handle carefully
\newcommand{\correction}{\textbf{[Point of correction.]}\ }
\newcommand{\concede}{\textbf{[Conceded.]}\ }
\newcommand{\newresult}{\textbf{[New result.]}\ }
\newcommand{\open}{\textbf{[Still open.]}\ }

% numbers that are not in yet, so they are impossible to miss on a read through
\newcommand{\tbd}[1]{\textbf{[TO FILL: #1]}}

% Marks what has been done since the first draft of these letters. The letters
% were written when most of this was a plan; this flag marks the places where it
% is now built and running.
\newcommand{\thisround}{\textbf{[Built since the first draft of this letter.]}\ }

\newcommand{\lettertitle}[2]{%
  \begin{center}
  {\Large\bfseries Housing First or Treatment First?}\\[0.3em]
  {\normalsize Manuscript RED-D-26-00078}\\[0.9em]
  {\large\bfseries #1}\\[0.3em]
  {\normalsize #2}
  \end{center}
  \vspace{0.6em}}


\begin{document}

\lettertitle{Letter to the Coordinating Editor}{Summary of changes}

Dear Professor Karabarbounis,

Thank you for your decision letter, and for leaving the door open to an almost new
paper. That is what I am resubmitting. The model, the calibration, the validation and
the experiments have all been rebuilt, and the paper has a new title. This letter
summarises what changed and which comments each change answers. Four further letters
respond to the Associate Editor and to each referee comment by comment, with the
responses in red. Section, table and figure numbers refer to the new paper. Comments are
numbered as in those letters: A1 to A20 for the Associate Editor, and R1, R2 and R3 for the
referees.

%======================================================================
\section*{\color{respred}What the paper now finds}

{\color{respred}
The paper is now \emph{Housing First or Treatment First? Ask the Kids, Not the
Parents}, and the title is its answer. The version you rejected found that, at the same
generosity, Treatment First reduced the disorder but raised homelessness, and it ranked
universal Treatment First above universal Housing First only once a loss of consortium
adjustment was added to welfare after the model was solved. The new paper drops the loss
of consortium and asks whether moral hazard, the argument that dominates the policy
debate, justifies making housing assistance conditional on sobriety.

It does not. Without any effect on children, a newborn is worse off under Treatment First
by 1.09 percent of lifetime consumption, even though drug use falls and the drug test's
savings buy vouchers for other households. Conditionality takes an option away from
forward looking households who weigh the risk of a substance use disorder (SUD)
correctly, and gives them back less through the government's budget than it takes.

The answer turns instead on the children, through three externalities that a parent
counts only in part. A child shares its parent's consumption and dwelling but not the
drug, so what the parent spends on drugs, and the earnings it loses while a SUD keeps it
out of work, the child goes without: the budget externality. A child whose parent had a
SUD while it lived at home starts adult life with lower productivity: the earnings
externality. And a child lives wherever its parent lives, so a revoked voucher can put it
in a shelter: the childhood homelessness externality, the one critics of Treatment First
expect it to create. Within the existing rationed programme, with the drug price and the
probability of drawing a voucher both clearing, Treatment First cuts drug use by 21
percent and the share of households with a SUD by 19 percent, and it lowers homelessness
and the share of children in a shelter by 11 percent each. A newborn gains 2.57 percent of
lifetime consumption, three quarters of it an externality its parent does not count. The
parents gain too, by 0.69 percent of their own lifetime consumption, because the higher
chance of drawing a voucher and the cheaper drug outweigh the option the screen removes.
Along the transition, neither homelessness nor the share of children in a shelter rises
in any year, and children born at the reform gain 1.45 percent.

The model reproduces the micro evidence. Without being targeted, it comes close to the
effect of making housing conditional on abstinence in the one trial that compares a sobriety
condition with unconditional housing (1.40 against 1.26; Milby et al., 2005), and in the
model losing a voucher raises the three year mortality of households with serious mental
illness (SMI) by about as much as Housing First lowered it among homeless veterans with
mental illness (3.6 against 4.6 percentage points; Costantini, 2025). As in that evidence,
households with a SUD lose: in the comparison a trial would make, with the drug price and the
voucher probability held, 74 percent of those with a SUD when the reform lands are worse off.
What reverses the verdict is what trials do not observe: the households the condition deters
from ever starting, the vouchers its savings buy for others, and the children.

The channel that decides the comparison came from the reports. Referee~2 (C3) asked for
parents who care about their children's prospects, and expected it to amplify the gain from
reducing the disorder. The Associate Editor (A6) objected that the previous version discussed
intergenerational transmission with no link between generations in the model. Built, the
channel does more than amplify: it reverses the ranking.
}

%======================================================================
\section*{\color{respred}What was built}

{\color{respred}
\textbf{A stylized model (Section~3)}, with two periods and calibrated (Tables~1 and~2),
shows why drug use concentrates at the bottom of the income distribution when no
preference varies with income (A1). Drugs are a luxury good, so without consequences the
average user would earn 3.51 times the population average. The risk of dying that a SUD
brings cannot be insured, and its cost grows with income; on its own it takes that figure to
0.64 at its calibrated strength. The risk of unemployment, which a household can partly insure by saving, leaves it at
3.51 even at twice its strength (Figure~1). Referee~1 asked for this comparison in the full
model (1(d)). I make it in a model with saving and with consequences that arrive later, the
two things the old static appendix lacked (1(e)), and in the quantitative model the mortality
multiplier of a SUD is calibrated to the income of users instead (Section~5.1).

\textbf{A quantitative model (Section~4)}, life cycle and annual, in general equilibrium in
the drug price and in the government's budget. Three things in it carry the paper.
\begin{itemize}
\item \emph{Children.} A child lives with its parent for 25 years, shares its consumption and
dwelling, and enters the parent's flow utility at the weight $\omega_c$. A SUD the parent has
while the child lives at home lowers the child's productivity for life, and the parent is
charged for that loss as a difference against a child whose parent never had the disorder
(Sections~4.2, 4.8 and 4.10).
\item \emph{Preferences that are the same in every state.} No preference parameter varies with
the disorder, with SMI or with housing (Section~4.2), so reducing the disorder or housing more
people raises no one's utility mechanically, which was the objection in Referee~2's C5,
Referee~3's M2(a) and the Associate Editor's A8. The disorder works through what it costs:
mortality, employment, recovery and the child's productivity. A child inherits a worse
productivity draw, not a taste for drugs.
\item \emph{Recovery that depends on housing.} Referee~2 (C1), the Associate Editor (A13) and
Referee~1 (2(d)) asked for discipline on the response of drug use to housing access. In the
model it runs through recovery, which is less likely in a shelter by a factor calibrated to
Milby et al.'s (2005) finding that housing raises the chance of recovering by 47.3 percent
(the model gives a ratio of 1.45 against 1.47). It is a transition probability rather than a
taste, and it is what the trial measures.
\end{itemize}
}

%======================================================================
\section*{\color{respred}What was cut, and what was added}

{\color{respred}
The Associate Editor asked whether a list of ingredients earned their place. I cut or
simplified each of them but one, the link between generations, which the paper now models and
which carries the result. Earlier versions solved the model with endogenous labour supply,
rents, homeownership and house prices, and these margins had a negligible effect on the
results (footnote~19).

\begin{center}
\small
\begin{tabular}{@{}L{3.7cm}L{9.0cm}L{2.2cm}@{}}
\toprule
Ingredient & Now & Comments \\
\midrule
Homeownership, mortgages and house prices & Cut. Households rent. The grid is the shelter and
three rental sizes; the shelter and the smallest rental are calibrated, and the two larger
sizes are set at four and eight times the smallest (Table~A1). & A4; R2 D1 \\[0.3em]
Endogenous rents & Cut. Equilibrium is now where the Associate Editor asked for it, in the
government's budget and the drug price, and Table~5 switches each on and off. & A5 \\[0.3em]
Labour supply & Cut. Labour is supplied inelastically when employed. & A3; R3 m1 \\[0.3em]
Unemployment insurance & One year of benefits for the share of job losers who are insured,
nothing after (Section~4.3). & A2 \\[0.3em]
Type dependent transfers and bequests & One endowment at entry, the same for every entrant,
from the estates of the dead; the government collects the rest. & A7 \\[0.3em]
Fiscal costs of homelessness and of a SUD & Cut. The shelter is free, and the budget counts
what the government pays in transfers, drug tests and treatment (Section~4.11). & A15;
R2 C6 \\[0.3em]
Light and heavy use & Cut. One chance that a year of use leads to a SUD (footnote~20). &
\\[0.3em]
Loss of consortium & Cut (below). & A9; R1 3(a) to 3(c) \\[0.3em]
Universal vouchers & Cut. Every experiment keeps the programme rationed, with the voucher
probability adjusting to hold net revenue at its benchmark (below). & R2 C6; R3 m3 \\
\bottomrule
\end{tabular}
\end{center}

Against that, the model gained what the reports asked for.
\begin{enumerate}
\item \textbf{SMI} (R1 1(a); Section~4.7): a persistent state whose onset a SUD makes twice as
likely, and which lowers the chance of finding a job, raises mortality and slows recovery from
a SUD. Without being targeted, households with SMI are 2.64 times as likely as others to use
drugs, against 2.44 in the data. Referee~1 expected it to weaken Treatment First; the screen
does fall hardest on households with SMI, and Treatment First still raises welfare.

\item \textbf{Treatment} (R1 1(b); A10; Section~4.6): a choice, at a utility cost calibrated to
a take up of 24.5 percent, that doubles the chance of recovering and that the government pays
for. Under Treatment First a household with a SUD passes the screen only in treatment and
without using, and the share in treatment rises from 24.6 to 29.0 percent. If a household with a SUD
cannot pass the screen at all, Treatment First still lowers homelessness by 11 percent and a newborn gains
at least as much (footnote~38).

\item \textbf{An equilibrium drug price} (R1 1(c); R3 M2(d); A5; Section~4.12), clearing a
supply curve with an elasticity of 1.46 (Mace et al., 2020). It works against Treatment First:
the price falls by 5 percent and draws users back, which takes the newborn's gain from 6.45 to
2.57 percent.

\item \textbf{Imperfect and costly enforcement} (R1 4(a) and 4(c); A11; Sections~4.9 and~6).
The screen catches a user with probability 20 percent a year, set so that it deters about as
much as the US military's random testing did: 62 percent of the users facing it stop, against
a fall of 61 percent in past year marijuana use (Mehay and Pacula, 1999). A test costs
35 dollars, the upper end of the estimates, and a household that is caught is excluded for
three years on average. At half the detection rate, Treatment First still cuts drug use by 14
percent and homelessness by 6 percent, and a newborn gains 1.65 percent.

\item \textbf{Transitions} (R1 3(d); Sections~6.4 and~6.5). The transition starts from the whole
Housing First population alive at the reform and clears the drug price and the voucher budget
every period (Figure~3). Drug use falls on impact, 81 percent of its long run fall in the first
year, and homelessness and the share of children in a shelter never rise; 11 percent of the
adults alive at the reform are worse off. When a household with a SUD cannot stop using
(Table~7), the design matters: a screen that catches every user from the first day raises
homelessness by up to 14 percent, while at the calibrated detection rate, exempting households
in treatment for the first five periods holds the rise to 1.6 percent and keeps the share of children in a shelter from rising
at all. Children born at the reform gain in every case.

\item \textbf{The disamenity of homelessness} (R2 C2). It is not in the benchmark. With one at
Loubi\`ere et al.'s (2020) contingent valuation, about 0.2 percent of average earnings a year, no
long run gain in Table~5 moves by more than 0.14 percentage points (footnote~4).

\item \textbf{The productivity cost of a SUD} (A20). The disorder lowers a user's earnings by
keeping it out of work longer (Section~4.8); there is no separate wage penalty on the user. The
productivity effect the model carries is on the child, set to Nilsson's (2017) natural
experiment estimate, 0.24. Even at zero, Treatment First raises welfare (footnote~5).
\end{enumerate}
}

%======================================================================
\section*{\color{respred}The loss of consortium adjustment is gone}

{\color{respred}
Referee~1 (3(a) to 3(c)), Referee~3 (M2(c)) and the Associate Editor (A9) all objected to it:
an adjustment applied after the model was solved, with no moment behind it, which changed the
ranking, and whose conservative value is arguably zero. I agree, and it is deleted. Welfare is
now a newborn's own expected lifetime utility, childhood included (Section~6.1), and nothing is
added to it after the model is solved.

What decides the ranking now lives inside the household's problem. A parent counts the share
$\omega_c$ of what its choices cost its child and acts on it; the rest is the externality. The
weight is calibrated to behaviour: in the 2024 National Survey on Drug Use and Health (NSDUH),
adults living with one other adult and no child are 17.5 percent more likely to use drugs than
comparable lone parents, and $\omega_c$ is the weight at which removing a parent's concern for
its child raises drug use in the model by as much (Section~5.1, Block~5). At Nishiyama's (2000) weight of 0.69, the upper end of the
literature, a newborn gains more, 3.18 percent, though less of it is an externality
(footnote~28). Section~6.3 then does what Referee~3 asked in M2(c): it splits the newborn's
gain into the childhood externality, 1.50 percent of lifetime consumption, the earnings
externality, 0.46, the newborn's own adult life under the policy, 0.43, and the part its parent
counts, 0.19 (Table~6).

On tone (Referee~1, 3(c)), the paper says who loses as plainly as who gains, and its
conclusion is conditional: yes if the externalities on children are counted, no if the case
rests on moral hazard (Section~7). In the long run 74 percent of households are better off and
3 percent worse off, but only 24 percent of those with a SUD are better off (Section~6.2).

\open The Associate Editor suggested reporting welfare and mortality separately, as a way round
the adjustment. With the adjustment gone, mortality enters welfare only through the value
households place on their own lives, which is calibrated to the value of a statistical life.
The paper checks the model's mortality response against Costantini (below), but it does not
report mortality as an outcome of the experiments.
}

%======================================================================
\section*{\color{respred}Errors in the previous version}

{\color{respred}
\concede Rebuilding the calibration from the data turned up errors in the version you
rejected. These are the ones that bore on the results.
\begin{enumerate}
\item \textbf{The homelessness target was a rate of people}, while the model's agent is a
household. The paper now targets a band of 0.40 to 1.00 percent of households, from the
Department of Housing and Urban Development's single night count to the share homeless at some
point in a year (Burt et al., 2001), and the model gives 0.62 percent.

\item \textbf{The share of the homeless with a SUD}, the 37 percent the Associate Editor
questioned (A16), was the midpoint of a wide range of substance abuse rates, not a rate of drug
use disorder. The paper now targets 21.7 percent, the drug use disorder share of the homeless
in Gutwinski et al. (2021), and the model gives 20.2. Homelessness means literal homelessness,
a shelter or the street (Section~4.5).

\item \textbf{The unemployment rate given a SUD was too high.} In the 2024 NSDUH it is 16.2
percent for adults aged 26 and over. It is no longer a target, and the model gives 14.5 percent
(Table~4).

\item \textbf{Drug use, the income of users and drug spending came from different definitions
of a user.} A user is now an adult aged 26 or over whose use amounts to a disorder of any
severity, 8.90 percent in the 2024 NSDUH. The income of users and the share of users with a
SUD are measured on that population, and drug spending is taken over its income
(Section~5.1, Block~1).
\end{enumerate}
}

%======================================================================
\section*{\color{respred}Calibration and validation}

{\color{respred}
Referee~1 could not assess the quantification because the calibration table never reported
the model's moments (2(a), 2(b)), and the Associate Editor agreed (A12). Of the model's 59
parameters, the 26 that carry the story are in Table~3, in six
blocks, and each entry says what the parameter controls, how it bears on the two policies and
what happens as it rises, before its value, moment and source (A17); the hazards of death and
unemployment are set out the same way, each with its source or moment (R3 M1(b)). The 33 standard ones are
in Appendix~A (Table~A1). Both tables report the data, the model and the value for every
target. The replication package documents the dataset, year, universe and computation behind
every estimate, and carries the data (A17; R1 2(e)).

Most targets are close. The largest gaps are the income of users, 62 percent of the average in
the model against 74 in the data, the share of renters paying more than half their income, 25.8
against 19.2 percent, and the unemployment rate, 5.6 against 4.9 percent. \concede The income of
users is the gap Referee~1 raised in 2(c)(i), where the previous model put it at 41 percent. He
was right that the model made users too poor, and it still does, by less. A new figure
(Figure~2) sets the share using drugs at each level of income against the 2024 NSDUH: use is
highest at the bottom and lowest at the top, as in the data, but in the model it is higher than
the data at the bottom and flatter in the middle, and both lower the income of users
(R3 M1(c) and M2(b)).

The policy responses the reports asked the model to match (R1 2(d); R2 C1; A13) are now in the
paper. Two are calibrated: the response of recovery to housing, to Milby et al. (Table~3), and
the response of drug use to a drug test, to the military's testing (Section~6). Two are not, and
the new validation section (Section~5.2, Table~4), which also compares the homeless,
users and voucher holders in the model with the data (R3 M1(a)), shows the model comes close to
both: Milby et
al.'s effect of making housing conditional on abstinence (1.40 against 1.26) and Costantini's
mortality effect (3.6 against 4.6). The price elasticities are no longer targets (R1 2(e); R3
M2(d)): participation responds more than in the evidence, $-1.32$ against $-0.3$ to $-0.9$,
which works against Treatment First, since the price it lowers draws users back. Transmission
across generations is at the low end of the evidence for a SUD, an odds ratio of 1.46 against
1.44 to 2.51, and above it for homelessness, 6.23 against 1.8 to 2.4, so the model loads the
effect on children towards the margin Housing First acts on.

The misses are reported as misses (R1 2(b), 2(c)). The homeless in the model are poorer and
younger than in the data: 21.4 percent are employed against 46, their income is 6.5 percent of
the average against 19.1, and 11.5 percent are 55 or older against 25.9. The paper reports
these as gaps and says what each means for the experiments.
}

%======================================================================
\section*{\color{respred}Costantini (2025)}

{\color{respred}
The Associate Editor (A14) read Costantini as finding that Housing First lowers the three
year mortality of homeless veterans by 4.6 percentage points while Treatment First has no
lasting effect on health, in contrast with the model. Validation Block~7 now asks the model
his Housing First question. With the drug price and the voucher probability held, losing a
voucher raises the three year mortality of the households with SMI whom Treatment First's
revocation puts on the street by about 3.6 percentage points, against his 4.6, without being
targeted. The paper reports only this arm. His Treatment First arm is the Department of
Veterans Affairs' own programmes, which the model does not represent, since its Treatment First
is a voucher whose retention depends on sobriety (Section~4.9), and I do not claim a
counterpart to it.

So the model agrees with him that losing housing raises mortality for the households who lose
it. What his design does not observe, and the model does, is everyone else: in the long run
Treatment First lowers homelessness by 11 percent, so fewer households are on the street at
all.
}

%======================================================================
\section*{\color{respred}Where I push back, and what the paper does not do}

{\color{respred}
Costantini is the first place, above. The others:
\begin{enumerate}
\item \textbf{Labour supply.} The Associate Editor (A3) asked what it was doing in the model,
and Referee~3 (m1) wanted it to price the cost of balancing the budget with labour taxes. I took
the first view. Taxes stay at their benchmark values in every experiment and the voucher
probability balances the budget, so no labour tax moves, and in earlier versions endogenous
labour supply had a negligible effect.

\item \textbf{The age limit on starting to use} (R1 S1). It stays, at age 80, and stops only
households without a SUD from starting. Without it, households close to the model's last age,
who know how little life remains, use too much (footnote~21).

\item \open \textbf{Generosity.} Referee~2 (C6) asked about optimal voucher design, Referee~3
(m3) for a continuous experiment in the voucher probability, and the Associate Editor found the
rationed and universal comparisons the most relevant. The paper now asks one question,
conditionality, within the existing rationed programme and at the same net cost. It does not
study how generous the programme should be.

\item \open \textbf{False positives} (R1 4(b)). The model has none, because a positive screen is
confirmed in a laboratory before it has consequences (Section~6), and there is no sensitivity to
them. Stigma and administrative burden (4(c)) are not modelled beyond the cost of a test, set
at the upper end of the estimates.

\item \open \textbf{Fiscal costs} (A15; R2 C6). Rather than defend the fiscal cost of
homelessness that both found high, I removed it, so the paper makes no comparison with Abramson
and Van Nieuwerburgh.

\item \open \textbf{Identification and sensitivity.} The paper reports neither the numerical
identification Referee~1 asked for (2(f)) nor a systematic sensitivity analysis (4(d)). It
reports the runs that bear on the ranking: half the detection rate, the parent's weight at the
upper end of the literature, no productivity penalty for the child, households with a SUD unable
to pass the screen, and a disamenity of homelessness. Treatment First raises a newborn's welfare
in each.
\end{enumerate}
}

%======================================================================
\section*{\color{respred}Presentation}

{\color{respred}
The body carries one comparison, Treatment First against Housing First in the rationed
programme, taken apart step by step in Table~5, decomposed in Table~6 and followed along its
transition in Figure~3 and Table~7. No experiments sit in an appendix, which holds only the
standard parameters (A18). The word ``addicted'' is gone (A19). Table~5 reports every outcome in
levels under each policy, and the text gives the changes in percent (R1 S2). The typographical
errors Referee~1 listed (S3 to S10) were in passages and references that have been rewritten or
removed. \open There are no plots of the drug use policy function (R2 C4) and no sensitivity
graphs (R3 m2); the paper shows who uses, and why, through Figures~1 and~2.
}

\vspace{1em}
Thank you again for the second look.

\vspace{1.4em}
Sincerely,\\[0.6em]
Shahar Rotberg

\end{document}
